\documentclass[a4paper,11pt]{article}
\usepackage[margin=27mm]{geometry}
\usepackage{fontspec,amsmath,amssymb,amsthm,unicode-math,booktabs,array}
\setmainfont{STIXGeneral}\setmathfont{STIX Two Math}
\setlength{\emergencystretch}{2em}
\newcommand{\gen}[1]{\langle#1\rangle}
\newtheorem{theorem}{Theorem}\newtheorem{proposition}[theorem]{Proposition}
\title{Rank intervals from one string marking}
\author{Human Compression project}\date{Private application draft; independent review pending}
\begin{document}\maketitle
\begin{abstract}
We apply the common contraction theorem to the principal results of
Brooksbank and Leemans, \emph{Rank reduction of string C-group
representations}. The same proof gives the stated rank intervals for
symmetric groups and for the retained classical-group constructions.
The highest-rank classical markings remain attributed inputs. We also
recover the source's finite matrix example and explain directly why
contraction of its three alternating-group diagrams loses generation.
No solution of the final rank-spectrum conjecture is asserted.
\end{abstract}

\section{The source theorem and its iteration}
The source is the complete six-page author edition arXiv:1812.01055v2,
3 May 2019, denoted by [BL]. Our common paper \emph{Contraction of string
markings through a central factor}, Theorem 3, proves its Theorem 1.1:
for an irreducible rank-$n$ string C-marking, $n\geq4$, the condition
\[
 r_0\in\gen{r_0r_2,r_3}
\]
ensures that $(r_1,r_0r_2,r_3,\ldots,r_{n-1})$ is a C-marking of the
same group in rank $n-1$. The common proof retains the standard end
criterion and proves the central-twist and dihedral steps in full.
It does not import [BL, Theorem 1.1] as an unproved premise.

If $r_2r_3$ has finite odd order $m$, then
$((r_0r_2)r_3)^m=r_0$. This recovers [BL, Corollary 1.2].
For iteration use the following precise form of common Corollary 5:
if $p_j=|r_{j-1}r_j|$ and $p_3,\ldots,p_{t+2}$ are odd, where
$0\leq t\leq n-3$, all ranks from $n-t$ through $n$ occur.
The first test is $p_3$, the next $p_4$, and so on. Irreducibility is
preserved, so each successive application has its full hypotheses.

The displayed maximum in the literal [BL, Corollary 1.3] additionally
tests $p_2$ and is undefined if $p_2$ is even. The version just stated
uses exactly the bonds required by its proof, permits zero steps without
any parity assumption, and implies every instance where that printed
maximum is defined. This indexing qualification does not change any
rank interval deduced below.

\section{Symmetric and alternating groups}
\begin{proposition}\label{symmetric}
For every $m\geq4$, the group $S_m$ has string C-markings of every
rank $3\leq n\leq m-1$.
\end{proposition}
\begin{proof}
Start with the adjacent transpositions
$(1\ 2),(2\ 3),\ldots,(m-1\ m)$. A selected set of these transpositions
generates the full symmetric group on each connected interval of the
selected path, independently on the different intervals. The intersections
of two such interval partitions are exactly the components of the
intersection of their selected edge sets: two vertices lie in the same
component precisely when every intervening edge is selected. Thus the
marking has the intersection property. Its adjacent products all have
order three and it generates $S_m$. Common Corollary 5 gives every stated
rank; $m=4$ is already the rank-three seed.
\end{proof}
This is the existence conclusion already proved by Fernandes and Leemans
in their 2011 symmetric-group paper, and explicitly deduced in [BL, \S3].
We make no claim here about uniqueness or enumeration of its highest-rank
representations.

For the alternating-group remarks in [BL, \S3], the necessary input is
more specific than the existence of some marking. We retain the highest-rank
markings of Fernandes--Leemans [FL], in the already acquired author edition
arXiv:1810.12450. For $m\equiv1\pmod4$, $m\geq13$, their Theorem 4.6
supplies rank $R=(m-1)/2$ and type $[10,3^{R-2}]$. Every bond tested by
contraction is three, so common Corollary 5 yields every rank $3$ through $R$.
For even $m\geq14$, the highest-rank instance underlying their Theorem 4.1
has rank $R=(m-2)/2$ and type
\[
 [5,6,3^{R-6},6,6,3].
\]
The consecutive odd bonds beginning at the third position give $R-6$
contractions, hence every rank from six to $R$. This recovers the
rank-existence observation made in [BL] about that theorem. The full
alternating-group classification in other congruence classes, its exact
lower-rank diagrams, and the nonexistence of particular ranks do not follow
from this one seed; their original constructions and classifications remain
outside this deduction.

\section{Classical groups: the construction inputs are retained}
We first specify the complete input used for characteristic two.
Brooksbank--Ferrara--Leemans [BFL, Corollary 1.5 and \S5] construct, for
$q=2^k$ with $k\geq2$ and $m\geq2$, string C-markings of
$O^+(2m,q)$ and $O^-(2m,q)$ in rank $2m$, and of
$\operatorname{Sp}(2m,q)$ in rank $2m+1$. The adjacent products in these
markings all have order $q+1$.

This input includes generation, the intersection property, and the
choice of Witt type. In the exact primary edition arXiv:1709.02219v2,
10 March 2018, the coefficients in \S5 are chosen so that the two-dimensional
products have order $q+1$; Lemma 5.3 and Proposition 3.3 supply generation
and intersections, and the proof of Corollary 1.5 chooses each Witt type
and a nonsingular radical in odd dimension. None of those construction
arguments is reclassified as a consequence of contraction.

\begin{theorem}\label{evenfield}
For integers $k,m\geq2$, $\operatorname{Sp}(2m,2^k)$ has string
C-markings in every rank $3\leq n\leq2m+1$, and each
$O^{\pm}(2m,2^k)$ has such markings in every rank $3\leq n\leq2m$.
\end{theorem}
\begin{proof}
Each retained highest-rank marking is irreducible: its adjacent order is
$q+1\geq5$. That order is odd. Apply common Corollary 5 until rank three.
The group is preserved at each contraction. This proves both clauses of
[BL, Theorem 1.4] with exactly the stated field and dimension restrictions.
\end{proof}
No assertion for $q=2$, no transfer to the derived orthogonal subgroup,
and no claim of maximality of these rank intervals is being made.

For the odd-characteristic result we retain the reflection construction
of Monson--Schulte, \emph{Reflection groups and polytopes over finite
fields, III}, \S3 and Theorem 3.1, together with its orthogonal-group
identification, as invoked in [BL, \S4]. The required form of that input is:
for $p\geq5$ prime and $d\geq3$, reduction of the crystallographic string
$[3,\infty^{d-2}]$ gives a rank-$d$ C-marking of an index-two subgroup
$M\leq O(d,p)$, with adjacent orders $[3,p^{d-2}]$.
The original Monson--Schulte PDF is not an acquired local primary in this
packet; its construction and group-identification theorem are explicitly
retained through the source's statement. This is a proof dependency, not
newly certified computation.

There is no hidden degeneracy in this chosen family. The integral Gram
matrix has diagonal entries $2$, first adjacent entries $-1$, remaining
adjacent entries $-2$, and zero elsewhere. Its leading determinants satisfy
$D_2=3$, $D_3=-2$, $D_4=-16$ and
\[
 D_j=2D_{j-1}-4D_{j-2}\quad(j\geq3),\qquad D_{j+3}=-8D_j\quad(j\geq2).
\]
The second identity follows from the first or by induction from the three
initial values. Every $D_d$ has no prime divisor other than two and three.
Consequently the form is nondegenerate after reduction at every $p\geq5$.
This calculation verifies the form condition; it does not replace the
retained generation or intersection theorem.

\begin{theorem}\label{oddfield}
For every $d\geq3$ and prime $p\geq5$, an index-two subgroup
$M\leq O(d,p)$ has string C-markings in every rank $3\leq n\leq d$.
\end{theorem}
\begin{proof}
All adjacent orders of the retained rank-$d$ marking are odd and greater
than two. Common Corollary 5 preserves its group $M$ and provides every
rank down to three. For $d=3$ the seed itself is the required marking.
This is [BL, Theorem 4.1]. The orthogonal type and the index-two subgroup
are those of the retained reflection construction; no identification with
an arbitrary prescribed simple projective group is inferred.
\end{proof}

\section{The matrix example outside the recovery hypothesis}
The membership condition in common Theorem 3 is sufficient for generation,
but it is not necessary for every successful contraction. We retain the
literal example of [BL, end of \S2] and verify it with the complete program
\texttt{check\_examples.g}. Over $\mathbb F_3$ put
\[
 F=\begin{pmatrix}1&1&0&0\\1&2&1&0\\0&1&1&2\\0&0&2&1\end{pmatrix},
\]
and use row-vector action with the following matrices:
\[
\begin{alignedat}{2}
 r_0&=\begin{pmatrix}2&0&0&0\\1&1&0&0\\0&0&1&0\\0&0&0&1\end{pmatrix},&
 r_1&=\begin{pmatrix}2&1&0&0\\0&1&0&0\\0&1&2&0\\0&0&0&2\end{pmatrix},\\
 r_2&=\begin{pmatrix}1&0&0&0\\0&1&1&0\\0&0&2&0\\0&0&2&1\end{pmatrix},&
 r_3&=\begin{pmatrix}2&0&0&0\\0&2&0&0\\0&0&2&1\\0&0&0&1\end{pmatrix}.
\end{alignedat}
\]
Direct multiplication gives $r_iFr_i^{\mathsf T}=F$. There are exactly
$21$ singular vectors, including zero, so this nonsingular four-dimensional
form has minus type. The generated matrix group has order $1440$, equal
to $|O^-(4,3)|=2\cdot3^2(3^2+1)(3^2-1)$; hence it is the full isometry
group, conventionally identified with $\operatorname{SL}(2,9):2$.
The last identification and the orthogonal order formula are standard
classical-group facts, not needed for the intersection check.

The program verifies all $136$ unordered pairs of parabolics of the
four-mark tuple and all $36$ for $(r_1,r_0r_2,r_3)$, including identical
subgroups and the trivial subgroup. It checks the types $[4,4,6]$ and
$[6,6]$, equality of the two generated groups, and
$r_0\notin\gen{r_0r_2,r_3}$. This is a complete finite verification of
the displayed example, using trusted GAP group computation; it is not
an argument from samples for common Theorem 3.

\section{The three $A_{11}$ diagrams and the remaining question}
Number the eleven vertices of each path in [BL, Figure 1] from left to
right. Its ten successive edges have the following label sets; a pair in
one entry denotes the displayed double edge.
\[
\begin{array}{c|llllllllll}
 &1&2&3&4&5&6&7&8&9&10\\\hline
\mathrm I&\{0,2\}&1&0&1&2&3&4&5&4&\{3,5\}\\
\mathrm {II}&0&1&\{0,2\}&1&2&3&4&5&4&\{3,5\}\\
\mathrm {III}&0&1&\{0,2\}&1&2&3&4&\{3,5\}&4&5
\end{array}
\]
The involution $r_i$ is the product of the path transpositions carrying
label $i$. Our complete program checks that each tuple generates $A_{11}$,
that all $2080$ unordered parabolic pairs satisfy IP, and that the three
types are $[5,3,6,3,5]$, $[5,5,6,3,5]$, and $[5,5,6,5,5]$.
Their different types, even after reversal, distinguish the three markings.
The statement that these exhaust the rank-six cases, and that $A_{11}$
has rank spectrum $\{3,6\}$, remains the prior classification cited in
[BL, \S3], Fernandes--Leemans--Mixer [FLM]. The finite check does not
prove exhaustion or the absence of ranks four and five.

For left contraction, the common transposition carrying labels $0$ and $2$
cancels in $r_0r_2$. It is the first path edge in row I and the third in
rows II and III. No other generator crosses that edge, so the new subgroup
is intransitive, with orbit sizes $(1,10)$ or $(3,8)$. From the right,
the analogous common edge of labels $3$ and $5$ cancels at the tenth edge
in rows I and II and at the eighth in row III. The corresponding orbit
sizes are $(10,1)$ or $(8,3)$. Thus neither direction generates $A_{11}$.
This proves the specific limitation in [BL] directly from all three
literal diagrams.

Conjecture 5.1 of [BL] asks whether $A_{11}$ is the only finite simple group
whose C-marking rank spectrum has a gap. Neither the contraction theorem
nor these examples settle that question. We retain the distinction between
an obstruction for a specified marking and an existence statement over
all markings of a group.

\section{Proof and computational boundaries}
The complete common theorem, its parity calculation and the applications
above recover [BL, Theorem 1.1, Corollaries 1.2--1.3, Theorems 1.4 and 4.1]
under the explicitly retained highest-rank construction inputs. The
iteration statement uses the precise indexing explained in Section 1.
The symmetric and selected alternating consequences and both families of
finite examples are also included. We do not claim to reprove the retained
classical seeds, the alternating classification, or the rank-gap conjecture.

The one whole GAP program is required for the two finite-example certificates.
The uniform proofs do not depend on its output. The four declarations in \texttt{RankContraction.lean} prove the odd-power
recovery identity, membership of the recovered generator, recovery of its
other factor, and exclusion from a generated subgroup by a homomorphism.
They are checked with Lean 4.32.2 and the recorded Mathlib revision. They do
not formalize the standard end criterion, the induction proving the C-property,
the classical-group seeds or the complete source-level application. No compression factor or
historical-priority claim is attached to this draft.

\begin{thebibliography}{9}
\bibitem{BL} P. A. Brooksbank and D. Leemans, \emph{Rank reduction of string
C-group representations}, arXiv:1812.01055v2, 3 May 2019;
Proc. Amer. Math. Soc. 147 (2019), 5421--5426.
\bibitem{BFL} P. A. Brooksbank, J. T. Ferrara and D. Leemans,
\emph{Orthogonal groups in characteristic 2 acting on polytopes of high rank},
arXiv:1709.02219v2, 10 March 2018; Discrete Comput. Geom. 63 (2020), 656--669.
\bibitem{FL} M. E. Fernandes and D. Leemans, \emph{String C-group
representations of alternating groups}, arXiv:1810.12450, exact acquired
author edition; Ars Math. Contemp. 17 (2019), 291--310.
\bibitem{FLM} M. E. Fernandes, D. Leemans and M. Mixer, \emph{Polytopes
of high rank for the alternating groups}, J. Combin. Theory Ser. A 119
(2012), 42--56.
\bibitem{MS} B. Monson and E. Schulte, \emph{Reflection groups and polytopes
over finite fields, III}, Adv. in Appl. Math. 41 (2008), 76--94.
\end{thebibliography}
\end{document}
